\chapter[Cryptography: Modular to Quantum]{Cryptography: Modular Integers to the Quantum Threat}

A student submits an assignment containing a cipher of his own making. Each letter is replaced by
its number, that number is multiplied by three, and the remainder after division by twenty-six is
used as the new letter. To open the message, he multiplies again by nine. When asked where the
number nine came from, his answer is nearly right: ``because there is no division here, I multiply
by the inverse.'' The examiner then replaces the multiplier three with the multiplier four. The
assignment is never finished. With the multiplier four, the first letter and the fourteenth letter
produce the same symbol, so two different messages produce an identical cipher.

The question that arises from that failure is the question of this whole chapter: when may a number
be divided inside remainder arithmetic? The answer is short: exactly when that number shares no
factor with its modulus. That short answer, pressed hard, yields the three things every browser
uses: key exchange, public keys, and digital signatures.

The previous three chapters are used explicitly. Chapter 2 gives the language for bijections and
partitions; here a bijection proves that multiplication by a unit only shuffles the order, and
partitions count the order. Chapter 3 gives the completeness property of the real numbers, and that
property does not cross over, because $\mathbb{Z}_n$ has no order compatible with multiplication;
what replaces it is finiteness, since on a finite set an injective map is automatically surjective.
Chapter 28 explains why ``secure'' cannot be a theorem without conditions.

One warning applies to every page here. All the moduli are small, chosen so that they can be checked
by hand and by Python in a matter of seconds. The largest modulus, $1\,050\,589$, is factored by a
six-line script. Not one cipher here is safe to use.

\begin{untukanak}
Imagine the wall clock in a classroom. Its hand moves from the number one to twelve, then back to
one. That clock never shows the number thirteen or twenty. If we turn the hand five steps, then
five steps again, then five steps again, we do not need to know how much five plus five plus five
is; we only look at where the hand stops.

Now imagine a combination lock with thirteen positions. We may turn it as far as we like, and only
the final position is recorded. If we turn three steps, how many times must we turn three steps
again to get back to where we started? The answer is nine, and that is not a coincidence, but
because three and thirteen share no factor. Try the same thing with four steps on a twelve-position
lock: however many times we turn, we will never land on all positions, because four and twelve
share a factor.

In this chapter we will understand why that small difference decides whether a message can be opened
again or is lost forever, and why two people who have never met can agree on one secret number in
front of a third person who hears everything.
\end{untukanak}

\section{Residue classes: arithmetic that goes around}
\tingkat{basic}

We begin by fixing the notion of ``remainder'' already used in the first paragraph.
\index{modular arithmetic}
Two integers are called congruent modulo $n$ if their difference is divisible by $n$.

\begin{definisi}
Let $n\ge1$ be an integer. For $a,b\in\mathbb{Z}$ we write $a\equiv b\pmod{n}$ if $n$ divides
$a-b$. The set of equivalence classes under this relation is called $\mathbb{Z}_n$, and the class
containing $a$ is written $[a]_n$.
\end{definisi}

We owe two things: that this really is an equivalence relation, and that the remainder is unique.

\begin{lemma}
Congruence modulo $n$ is an equivalence relation on $\mathbb{Z}$, and every class $[a]_n$ contains
exactly one representative in $\{0,1,\dots,n-1\}$, so that $|\mathbb{Z}_n|=n$.
\end{lemma}

\begin{proof}
Reflexive: $n$ divides $0$. Symmetric: if $n\mid(a-b)$ then $n\mid(b-a)$. Transitive: if
$n\mid(a-b)$ and $n\mid(b-c)$ then $n\mid(a-c)$ because $(a-c)=(a-b)+(b-c)$. For uniqueness, take
$a$ and consider $S=\{a-qn\ge0: q\in\mathbb{Z}\}$. That set is nonempty by the Archimedean property
(Chapter 3), so it has a least element $r$. If $r\ge n$ then $r-n\in S$ is smaller than $r$, a
contradiction; so $0\le r<n$. If $r'$ also satisfies $0\le r'<n$, then $r-r'$ is a multiple of $n$
whose value lies between $-(n-1)$ and $n-1$, hence must be zero.
\end{proof}

Operations on $\mathbb{Z}_n$ are defined through their representatives, and a definition of this
kind is dangerous: if the result depends on the representative, the operation is not a function. The
next lemma closes that hole with one line of algebra.

\begin{lemma}
If $a\equiv a'$ and $b\equiv b'$ modulo $n$, then $a+b\equiv a'+b'$ and $ab\equiv a'b'$ modulo $n$.
\end{lemma}

\begin{proof}
Write $a'=a+kn$ and $b'=b+ln$. Then $a'+b'=(a+b)+(k+l)n$ and $a'b'=ab+(al+bk+kln)n$.
\end{proof}

Addition and multiplication are therefore well defined, and both inherit associativity and
distributivity from $\mathbb{Z}$. What is lost is order: there is no way to put a sign ``$<$'' on
$\mathbb{Z}_n$ such that $a<b$ implies $a+c<b+c$ and $ac<bc$ for positive $c$. The supremum tool of
Chapter 3 stops at the entrance, and what works inside is counting.

Division need not exist. The class $[3]_{26}$ has an inverse, namely $[9]_{26}$, since
$3\cdot9=27\equiv1$. The class $[2]_{26}$ has none: every multiple of two is even, whereas $27$ is
odd, so $2b\equiv1$ has no solution. Consequently multiplication by two modulo twenty-six is not a
bijection, because $[0]$ and $[13]$ both go to $[0]$.

\kamus{residue class}{kelas sisa}{all the numbers that stop at the same clock position, like three
in the morning and three in the afternoon}

\begin{definisi}
A class $[a]_n$ is called a unit if there is $[b]_n$ with $[a]_n[b]_n=[1]_n$. The set of all units
modulo $n$ is written $U(n)$, and its number of members is written $\varphi(n)$.
\end{definisi}

For $n=26$ there are twelve units, namely the classes of $1,3,5,7,9,11,15,17,19,21,23,25$. The
program in the lab counts them again; that number reappears in part three under a different name.

\section{Common factors, Euclid, and when division exists}
\tingkat{basic}

\index{greatest common divisor}
The question ``when does $[a]_n$ have an inverse'' turns out to be a question about the greatest
common factor, which has an algorithm more than two thousand years old.

\begin{definisi}
For $a,b\in\mathbb{Z}$ not both zero, $\gcd(a,b)$ is the largest natural number that divides $a$ and
$b$.
\end{definisi}

The key to Euclid's algorithm is one equality of sets: the common divisors of $a$ and $b$ are exactly
the common divisors of $b$ and $a-qb$.

\begin{lemma}
For every $q\in\mathbb{Z}$, $\gcd(a,b)=\gcd(b,a-qb)$.
\end{lemma}

\begin{proof}
If $d$ divides $a$ and $b$, then $d$ divides $a-qb$. Conversely, if $d$ divides $b$ and $a-qb$, then
$d$ divides $(a-qb)+qb=a$. So the two pairs have the same set of common divisors, hence the same
greatest one.
\end{proof}

Replacing $(a,b)$ by $(b,a\bmod b)$ therefore does not change the answer but shrinks $b$, and the
process stops because natural numbers do not shrink without bound.

\begin{teorema}[Bezout]
For $a,b$ not both zero, there are $x,y\in\mathbb{Z}$ with $ax+by=\gcd(a,b)$.
\end{teorema}

\begin{proof}
Consider $S=\{ax+by>0: x,y\in\mathbb{Z}\}$. This set is nonempty because it contains $|a|$ or $|b|$,
and every nonempty subset of the natural numbers has a least element; this is the discrete form of
the infimum property in Chapter 3. Let $d=ax_0+by_0$ be that least element. Divide $a$ by $d$:
$a=qd+r$ with $0\le r<d$. Then $r=a(1-qx_0)+b(-qy_0)$, so $r\in S$ if $r>0$. Since $d$ is least,
$r>0$ is impossible; so $r=0$ and $d\mid a$, and in the same way $d\mid b$. Hence $d\le\gcd(a,b)$.
Conversely $\gcd(a,b)$ divides every expression $ax+by$, in particular $d$, so that
$\gcd(a,b)\le d$.
\end{proof}

\begin{teorema}
The class $[a]_n$ is a unit if and only if $\gcd(a,n)=1$.
\end{teorema}

\begin{proof}
If $[a]_n$ is a unit, there is $b$ with $ab\equiv1$, that is $ab-1=kn$. Every common divisor of $a$
and $n$ divides $ab$ and $kn$, hence divides $1$; so $\gcd(a,n)=1$. Conversely, if $\gcd(a,n)=1$,
Bezout's theorem gives $ax+ny=1$. Read modulo $n$, that equality says $ax\equiv1$, so $[x]_n$ is the
inverse of $[a]_n$.
\end{proof}

That proof also gives a recipe: run Euclid, then build the combination backward. For $a=1071$ and
$b=462$ the steps are three divisions,
\[
1071=2\cdot462+147,\qquad 462=3\cdot147+21,\qquad 147=7\cdot21+0,
\]
so that $\gcd=21$. Building backward gives $21=1071\cdot(-3)+462\cdot7$, and since the common factor
is $21$ rather than $1$, the class $[462]_{1071}$ is not a unit.

How long does this algorithm run? Every two steps the remainder at least halves, so the number of
divisions is proportional to the number of digits. The sharp bound was given by Lam\'e in 1844: no
more than $\log_\phi(b)+1$ divisions with $\phi=(1+\sqrt5)/2$, and the pairs that attain it are
consecutive Fibonacci pairs. The pair $(F_{47},F_{46})=(2\,971\,215\,073,\ 1\,836\,311\,903)$ needs
$45$ divisions, whereas $\log_\phi(F_{46})=44.33$. For two thousand random thirty-two-bit pairs, the
average is $18.75$ with maximum $32$.

\kamus{modular inverse}{invers modular}{the number that, multiplied by the original number, gives
remainder one, like the ``undo'' button on a cash register}

\section{Fermat's little theorem and Euler}
\tingkat{intermediate}

This chapter needs one fact about powers, and that fact is born from counting: multiply every unit by
one fixed unit; the result is the same list in a different order.

\begin{lemma}
If $\gcd(a,n)=1$, the map $[x]_n\mapsto[ax]_n$ is a bijection from $U(n)$ to $U(n)$.
\end{lemma}

\begin{proof}
The result lies in $U(n)$ because the product of two units is a unit: if $xx'\equiv1$ and
$yy'\equiv1$ then $(xy)(x'y')\equiv1$. The map is injective: if $ax\equiv ay$, multiply both sides by
the inverse of $a$, which exists by the unit theorem in part two. The set is finite, and on a finite
set an injective map from the set to itself is automatically surjective (Chapter 2). So it is a
bijection.
\end{proof}

\begin{teorema}[Euler]
If $\gcd(a,n)=1$, then $a^{\varphi(n)}\equiv1\pmod n$.
\end{teorema}

\begin{proof}
Let $u_1,\dots,u_k$ be all the units modulo $n$, with $k=\varphi(n)$, and $P=u_1\cdots u_k$. By the
lemma above, the list $au_1,\dots,au_k$ contains every unit exactly once, so
\[
a^kP\equiv(au_1)(au_2)\cdots(au_k)\equiv u_1u_2\cdots u_k=P\pmod n .
\]
Since $P$ is itself a unit, it has an inverse, and both sides may be multiplied by it. The result is
$a^k\equiv1$.
\end{proof}

\begin{akibat}[Fermat's little]
If $p$ is prime and $p\nmid a$, then $a^{p-1}\equiv1\pmod p$. For all $a$ we have
$a^p\equiv a\pmod p$.
\end{akibat}

The second form follows from the first: if $p\mid a$ both sides are zero, if not, multiply
$a^{p-1}\equiv1$ by $a$. The name ``little'' distinguishes it from Fermat's last theorem, which is
far harder.

The number twelve from part one now has two names: $\varphi(26)=12$, and $3^{12}\equiv1\pmod{26}$;
the program checks that $5^{12}$ and $7^{12}$ are also $\equiv1$. Large powers can therefore be
trimmed: for $7^{222}\bmod13$ write $222=12\cdot18+6$, so that $7^{222}\equiv7^6\equiv9\cdot10\equiv12$
because $7^2\equiv10$ and $7^4\equiv9$.

The converse does not hold, and that matters for the honesty of this chapter. The number
$561=3\cdot11\cdot17$ is not prime, yet $a^{560}\equiv1\pmod{561}$ for every $a$ that shares no
factor with it; the program checks all $a$ without finding a counterexample. Numbers of this kind are
called Carmichael, and for that reason practical primality tests use Miller--Rabin.

\kamus{Euler totient function}{fungsi phi}{the number of positions on a combination lock that can
reach every other position, that is the positions that share no factor with the number of positions}

\section{Cyclic groups, order, and primitive roots}
\tingkat{intermediate}

If we multiply one unit repeatedly, the result goes around and finally returns to one. The useful
question is how long that circuit is.

\begin{definisi}
For $\gcd(a,n)=1$, the order of $a$ modulo $n$ is the least natural number $k$ with
$a^k\equiv1\pmod n$, written $\operatorname{ord}_n(a)$.
\end{definisi}

An order always exists: the sequence $a^0,a^1,\dots$ has only $\varphi(n)$ distinct values, so two of
them are equal and $a^{i-j}\equiv1$ for some $i>j$. That bound can be sharpened: the order divides
$\varphi(n)$.

\begin{lemma}
If $\gcd(a,n)=1$, then $\operatorname{ord}_n(a)$ divides $\varphi(n)$.
\end{lemma}

\begin{proof}
Let $k=\operatorname{ord}_n(a)$ and $H=\{a^0,\dots,a^{k-1}\}$, whose members are all distinct, so
$|H|=k$. For every unit $x$, the set $xH=\{xa^0,\dots,xa^{k-1}\}$ also has $k$ members, because
multiplication by a unit is a bijection. The sets $xH$ form a partition of $U(n)$: two of them are
equal or disjoint, since $xH$ and $yH$ intersect if and only if $y^{-1}x\in H$ (Chapter 2). So
$\varphi(n)$ is $k$ times the number of parts.
\end{proof}

\begin{definisi}
If $\operatorname{ord}_n(g)=\varphi(n)$, then $g$ is called a primitive root modulo $n$. For a unit
$h$, the number $k$ with $g^k\equiv h$ is called the discrete logarithm of $h$ to the base $g$.
\end{definisi}

Modulo $13$, $2,6,7,11$ have full order $12$ and are therefore primitive roots, whereas $3$ and $9$
have order $3$, $5$ and $8$ have order $4$, $4$ and $10$ have order $6$, and $12$ has order $2$. All
those numbers divide $12$.

For $2$, the circuit visits all twelve nonzero classes before returning: $2,4,8,3,6,12$, then
$11,9,5,10,7,1$. Figure~\ref{fig:jam} shows it as a clock with thirteen positions. For $3$, the
circuit is only $3,9,1$ and then repeats.

\begin{figure}[htbp]
\centering
\begin{tikzpicture}
\foreach \k in {0,...,12}
  \node[circle,draw,inner sep=1pt,minimum size=6mm,font=\scriptsize] (n\k)
    at ({90-360*\k/13}:2.1cm) {\k};
\node[font=\footnotesize,align=center] at (0,0) {multiply\\ by 2};
\foreach \a/\b in {1/2,2/4,4/8,8/3,3/6,6/12,12/11,11/9,9/5,5/10,10/7,7/1}
  \draw[-{Latex[length=1.5mm]},semithick] (n\a) to[bend left=18] (n\b);
\end{tikzpicture}
\caption{Multiplication by $2$ modulo $13$ goes around through all twelve nonzero classes before
returning to $1$. The length of that circuit is the order, and full order means $2$ is a primitive
root modulo $13$.}
\label{fig:jam}
\end{figure}

The existence of a primitive root for every prime $p$ is Gauss's theorem, and we skip its proof. What
must be carried away: for a prime $p$, the units are a single circuit of $p-1$ positions, and every
unit is a power of one primitive root. If $h=g^k$, knowing $k$ means knowing $h$, but not
conversely; that asymmetry is the raw material of public-key cryptography.

\kamus{primitive root}{akar primitif}{one starting position that, if turned continuously, passes
through every other position before returning}

\section{Diffie-Hellman key exchange}
\tingkat{intermediate}

\index{diffie-hellman}
Two people want to agree on one secret number, and the only channel they have is overheard by a third
party. In 1976 Diffie and Hellman showed that this is possible, and the mechanism uses only two
formulas.

They first agree on a prime $p$ and a primitive root $g$ modulo $p$; both may be announced. Alice
chooses $a$ secretly and sends $A=g^a\bmod p$. Bob chooses $b$ secretly and sends $B=g^b\bmod p$.
Alice computes $K=B^a\bmod p$, Bob computes $K=A^b\bmod p$. Both obtain the same number, because
\[
B^a\equiv(g^b)^a\equiv g^{ab}\equiv(g^a)^b\equiv A^b\pmod p .
\]
What the eavesdropper sees is only $p,g,A,B$. To obtain $K$ he needs $g^{ab}$, and the easiest known
way is to find $a$ or $b$, that is to solve the discrete logarithm.

\begin{bacakeras}
Read it as a dialogue for two voices, alternating, one line per breath.

First voice: ``Agree first, announce two.''
Second voice: ``Keep your own, never say it.''
First voice: ``Raise, take the remainder, send the result.''
Second voice: ``Receive the delivery, raise it again.''
First voice: ``The last remainder, that is the key.''
Together: ``Two directions, one number, never crossed.''

Repeat three times, faster and faster. The line that makes the tongue stumble is the step not yet
understood, not the line that needs to be memorized harder. Once it flows, close the book and name
four things in order: what is announced, what is kept, what is sent, what is computed. Naming the
secret number first cancels the whole game.

The line repeated three times aloud: ``what is sent is the power, what is kept is the exponent.''
\end{bacakeras}

The concrete numbers, all checked by the program in the lab part: $p=10\,007$, $g=5$ has order
$10\,006$ and is therefore a primitive root, $a=1234$, $b=5678$. Then $A=5^{1234}\bmod p=4873$ and
$B=5^{5678}\bmod p=3189$. Alice computes $3189^{1234}\bmod p=5674$; Bob computes
$4873^{5678}\bmod p=5674$. Their shared key is $5674$, and that number never crosses the channel.
Figure~\ref{fig:dh} shows the flow together with what the eavesdropper hears.

\begin{figure}[htbp]
\centering
\begin{tikzpicture}[font=\footnotesize,
  kotak/.style={draw,rounded corners=2pt,align=left,inner sep=4pt,text width=4.4cm}]
\node[kotak] (a) {\textbf{Alice}\\ $a=1234$ secret\\ $A=5^{1234}\bmod p=4873$};
\node[kotak,right=3.4cm of a] (b) {\textbf{Bob}\\ $b=5678$ secret\\
  $B=5^{5678}\bmod p=3189$};
\node[kotak,below=1.3cm of $(a)!0.5!(b)$, text width=9.2cm, fill=black!4]
  (e) {\textbf{Eavesdropper} hears $p=10\,007$, $g=5$, $A=4873$, $B=3189$, then intends
  to compute $K=A^{b}$ or $K=B^{a}$; for that it must find $a$ or $b$.};
\draw[-{Latex}] ([yshift=3mm]a.east) -- node[above,font=\scriptsize]{$A$}
  ([yshift=3mm]b.west);
\draw[-{Latex}] ([yshift=-3mm]b.west) -- node[below,font=\scriptsize]{$B$}
  ([yshift=-3mm]a.east);
\draw[-{Latex},dashed] (e.north -| a.south) -- (a.south);
\draw[-{Latex},dashed] (e.north -| b.south) -- (b.south);
\end{tikzpicture}
\caption{The flow of Diffie-Hellman key exchange on the numerical example in the text. The values $A$
and $B$ cross the open channel; what does not cross are the exponents $a$ and $b$, and it is exactly
those two that the eavesdropper needs.}
\label{fig:dh}
\end{figure}

That equality only states that the protocol is consistent, not that it is secure. All that can be
proved is a conditional statement: if the discrete logarithm modulo $p$ is hard, then computing
$g^{ab}$ from $A$ and $B$ is also hard. That condition is not a theorem; there are groups in which
the discrete logarithm is easy, so security is a property of particular groups.

\begin{tangan}
Materials: yellow watercolor, one secret color per person (red and blue), four plastic cups, one
spoon, two sheets of white paper.

First step, in front of the table: pour one spoonful of yellow paint into each cup. The yellow plays
the role of the prime and the base, and both are announced. Second step, behind your back: each
person adds two spoonfuls of their secret color, stirs, then hands over the mixed cup. Third step:
each adds two spoonfuls of their own secret color to the cup just received, stirs, then puts one drop
on the paper.

What must be visible: the two drops are the same color, even though neither of the two people ever
handed over their secret color. The first mixed cup held by the third person cannot be separated
again; he sees only two different dull colors.
\end{tangan}

Two parameter traps deserve mention. First, $g$ must generate a large subgroup of prime order: the
modulus $10\,007$ is a safe prime with $p-1=2\cdot5003$, and the program records that
$5^{5003}\bmod p=10\,006$ has order only $2$. If Alice uses that base, their shared key is only $1$
or $10\,006$, and the eavesdropper need not compute any logarithm at all. Second, this protocol
names no one's identity, so an eavesdropper can stand in the middle and read everything.
Figure~\ref{fig:dh} depicts only a passive eavesdropper.

About size: the modulus $10\,007$ can be broken in less than one second by the index calculus
algorithm, and we use it because the powers fit on one page. Real practice uses $2048$ bits or more,
and even that is a time limit, not a guarantee.

\begin{untukanak}
Imagine two people who want to agree on one secret color, even though all their conversation is
overheard. They first agree on a base color, say yellow. Alice chooses her secret color, mixes it with
yellow, then sends the mixture. Bob does the same with his own secret color. Now each mixes the color
received with their own secret color. Both get the same color, and the eavesdropper sees only two
mixtures that cannot be separated again.

That separation is what is considered hard: mixing colors is easy, separating a mixture without one
of its ingredients is considered not. Cryptography replaces colors with numbers and ``considered not
easy'' with a mathematical sentence.

In this chapter we will find out how numbers replace colors, and why the security depends on one
``considered'' sentence that nobody has ever proved.
\end{untukanak}

\section{RSA: from key to message and back}
\tingkat{intermediate}

\index{rsa}
Diffie-Hellman agrees on a key, but does not send a message. A year later Rivest, Shamir, and
Adleman showed that the same property can be used for two different keys: one to lock, one to open.

Choose two large primes $p$ and $q$, compute $n=pq$ and $\varphi(n)=(p-1)(q-1)$. Choose $e$ with
$\gcd(e,\varphi(n))=1$; the unit theorem gives its inverse $d$ modulo $\varphi(n)$. The public key is
$(n,e)$, the secret key $d$. For a message $0\le m<n$, the cipher is $c=m^e\bmod n$ and its opening is
$m=c^d\bmod n$.

Why does opening return $m$? Because $ed=1+k\varphi(n)$, so that
$c^d=m^{ed}=m\,(m^{\varphi(n)})^k$. If $\gcd(m,n)=1$, Euler's theorem makes that bracket one. The
remaining case explains why $n$ must be the product of two distinct primes.

\begin{teorema}
Let $n=pq$ with $p,q$ distinct primes, $ed\equiv1\pmod{(p-1)(q-1)}$, and $0\le m<n$. Then
$m^{ed}\equiv m\pmod n$.
\end{teorema}

\begin{proof}
If $p\nmid m$, Fermat's little theorem gives $m^{p-1}\equiv1\pmod p$, so that
$m^{ed}=m\,(m^{p-1})^{k(q-1)}\equiv m\pmod p$. If $p\mid m$, both sides are congruent to zero modulo
$p$. So the congruence holds modulo $p$ for every $m$, and for the same reason modulo $q$. Since $p$
and $q$ both divide $m^{ed}-m$ and $\gcd(p,q)=1$, their product also divides that number: if
$m^{ed}-m=pt$ and $q\mid pt$ with $q$ prime and $q\nmid p$, then $q\mid t$, so that $pq\mid m^{ed}-m$.
\end{proof}

The numerical example is small and entirely checked by the program. Take $p=1019$, $q=1031$, then
$n=1\,050\,589$ and $\varphi(n)=1018\cdot1030=1\,048\,540$. Choose $e=65\,537$; Euclid gives
$d=60\,493$, and indeed $65\,537\cdot60\,493=3\,964\,529\,741=1+3781\cdot1\,048\,540$. The message
$m=424\,242$ produces the cipher $c=552\,241$, and $c^d\bmod n$ returns $424\,242$. That modulus is
only twenty-one bits and is factored by a six-line script.

Plain RSA, the formula above with no addition of any kind, is not secure; its three reasons can be
shown with the same modulus. First, it is deterministic: the same message always produces the same
cipher, so that repeated messages can be recognized. Second, it is multiplicative:
$E(m_1)E(m_2)\equiv E(m_1m_2)$, and the program shows
$E(1234)\cdot E(567)\bmod n=551\,173$ opened to $699\,678$, exactly $1234\cdot567$. Third, small
messages are not protected: with $e=3$ and $m=42$, the cipher is $74\,088$, coincidentally equal to
$42^3$ because the result has not yet passed $n$, so that an ordinary cube root recovers the message.

All three are closed by padding: the message receives a random addition before exponentiation, so
that the result is full length, always different, and no longer multiplicative. RSA-OAEP, the
standard used today, has a security proof against chosen-ciphertext attacks in the random oracle
model under the assumption that RSA holds. In practice RSA wraps a $256$-bit symmetric key and then
the data is encrypted with AES; this pattern is called hybrid encryption, because RSA is slow: its
time is proportional to the cube of the modulus length.

\kamus{padding}{bantalan}{extra contents put into a package before shipping, like bubble wrap that
makes every package a different shape and size}

\section{Hash functions, MACs, and signatures}
\tingkat{intermediate}

\index{hash function}
Encryption hides the contents. Sometimes what is needed is not hiding but summarizing: one short
fingerprint that changes when the contents change.

\begin{definisi}
A cryptographic hash function maps a message of any length to an output of fixed size. It must be
preimage resistant (from $h$, it is hard to find $m$ with $H(m)=h$), second-preimage resistant (from
$m$, it is hard to find $m'\neq m$ with $H(m')=H(m)$), and collision resistant (it is hard to find a
distinct pair $m,m'$ with the same hash).
\end{definisi}

Those three properties are not equivalent; the weakest is the one most often misunderstood. Collision
resistance is limited by birthdays: if the output has $N$ values, a collision appears after about
$1.177\sqrt N$ trials, not after $N$; for $128$ bits that number is $2.17\cdot10^{19}$ and for $64$
bits $5.06\cdot10^{9}$. The program in the lab tests that prediction on SHA-256 truncated to $16$
bits: an average of $362.7$ samples against a prediction of $320.8$, with a spread from $100$ to
$743$. So a hash with an $n$-bit output gives only $n/2$ bits of collision resistance.

A MAC adds a key to the hash: two parties who share a key can check that a message came from the
other party and did not change, because without the key there is no way to compute the correct tag.
The limit: a MAC is useful only between two parties who both hold the key, so that a third party
cannot be a witness.

A signature uses a key pair in the opposite direction: the signer uses the secret key, the verifier
uses the public key. In RSA, the signature on $m$ is $s=H(m)^d\bmod n$, and verification compares
$s^e\bmod n$ with $H(m)$. The need for $H$ is not a matter of convenience: without a hash, RSA
signatures are multiplicative like its encryption, so that from the signatures of two documents
someone can assemble a valid signature for their product without holding the secret key. The hash
breaks that property, because $H(m_1m_2)$ is unrelated to $H(m_1)H(m_2)$. The security requirement is
resistance to existential forgery under chosen-message attack, and once again that is an assumption.

\kamus{digital signature}{tanda tangan digital}{a wax seal on a document that anyone can check but
only the holder of the stamp can make}

\section{What the security model actually assumes}
\tingkat{advanced}

Every sentence with the word ``hard'' in this chapter is a promise about what the known algorithms
cannot do, not a statement about what is impossible.

\begin{definisi}
A function $f$ is called one-way if there is a polynomial algorithm to compute $f(x)$ from $x$, but
for random $y$ in the range, every polynomial algorithm fails to find $x$ with $f(x)=y$ except with
negligible probability.
\end{definisi}

The direction is depicted in Figure~\ref{fig:tangga-satu-arah}.

\begin{figure}[htbp]
\centering
\begin{tikzpicture}[font=\small,
  kotak/.style={draw,rounded corners=2pt,align=center,inner sep=5pt,
    minimum height=0.95cm}]
  \node[kotak,draw=layerbrown] (pq) at (0,0) {$p$, $q$ random primes};
  \node[kotak,draw=ink] (n) at (4.0,0) {$n=p\,q$};
  \node[kotak,draw=layergrey,dashed] (s) at (8.2,0) {attacker};
  \draw[-{Latex},very thick,layerjade] (pq) -- (n);
  \node[font=\scriptsize,layerjade,above] at (2.0,0.06)
    {multiply: easy, polynomial};
  \draw[-{Latex},thick,layergrey] (n) -- (s);
  \node[font=\scriptsize,layergrey,above] at (6.1,0.06) {sees $n$};
  \draw[layerplum,thick,dashed] (n.south) -- (4.0,-1.35);
  \draw[layerplum,thick,dashed] (0,-1.35) -- (pq.south);
  \draw[-{Latex},layerplum,thick,dashed] (4.0,-1.35) -- (0,-1.35);
  \draw[layerplum,very thick] (1.86,-1.12)--(1.86,-1.58);
  \draw[layerplum,very thick] (2.04,-1.12)--(2.04,-1.58);
  \node[font=\scriptsize,layerplum,below] at (2.0,-1.62)
    {factor: hard, no known polynomial algorithm};
\end{tikzpicture}
\caption{Two directions that are not equally expensive. From $p$ and $q$ to $n$ one multiplication
suffices, whereas from $n$ back to $p$ and $q$ is a problem that to this day has no known polynomial
algorithm; that gap is what RSA and Diffie-Hellman use. The bar on the return arrow marks a direction
assumed closed, not a direction already proved impossible.}
\label{fig:tangga-satu-arah}
\end{figure}

That definition does not name a single function, and that is not an oversight. To this day there is
no proof that one-way functions exist; their existence would imply $P\neq NP$, which is also still
open. Public-key cryptography therefore stands on a conjecture, and the chain ends in an assumption
that does not turn into a theorem by being used for a long time.

\index{factoring assumption}
The factoring assumption says: for $n$ the product of two random primes of size $k$ bits, there is no
polynomial algorithm in $k$ that factors it. The best known algorithm, the number field sieve, runs
in time $L_n[1/3,1.923]$, which grows subexponentially. The number RSA-250, $250$ digits long or
about $829$ bits, was factored in 2020 at a cost of about $2\,700$ core-years.

The strength equivalences recorded by NIST in SP 800-57 are useful as a rough map.

\begin{center}
\small
\begin{tabular}{@{}lll@{}}
\toprule
Strength (bits) & RSA or DH modulus (bits) & Symmetric cipher \\
\midrule
112 & 2048 & 3DES three keys \\
128 & 3072 & AES-128 \\
192 & 7680 & AES-192 \\
256 & 15360 & AES-256 \\
\bottomrule
\end{tabular}
\end{center}

The numbers in the middle column are cost estimates for the best algorithm at the time the table was
compiled, so a table of this kind has an expiry date.

Formal security models close the gap in the rhetoric with a game: the challenger generates keys,
answers encryption queries, then chooses two messages and sends the encryption of one of them; the
scheme is secure if the attacker cannot guess which one far above the chance of one half. Security is
proved as a reduction: an attacker who wins is turned into an algorithm that factors $n$, so that all
the doubt moves to one place that is narrow but does not disappear.

\section{Shor, and what collapses with it}
\tingkat{advanced}

\index{Shor's algorithm}
The whole building above depends on one difficulty, and in 1994 Peter Shor showed that a quantum
computer solves it in polynomial time.

The idea splits in two, and the first part is entirely classical. To factor $n$, choose $a$ with
$\gcd(a,n)=1$ and compute $r=\operatorname{ord}_n(a)$. If $r$ is odd, repeat with another $a$. If $r$
is even and $a^{r/2}\not\equiv-1\pmod n$, then $\gcd(a^{r/2}-1,n)$ is a proper factor of $n$: because
$a^r\equiv1$, the number $a^{r/2}$ is a square root of one that is not $\pm1$, and a root of that kind
appears only when the modulus splits. The rest of the work is finding the order, and this is the
quantum part.

The classical part alone is not enough. For $n=1019\cdot1031$ we tried $2\,996$ random units and only
$319$, about $10.6$ percent, satisfied the condition and so gave a factor; success therefore comes
from repeating with different $a$. What is missing without a quantum computer is a fast way to find
$r$: no known classical algorithm runs in polynomial time in $\log n$.

A quantum computer finds $r$ with the quantum Fourier transform, and that is where the speedup comes
from. The consequence is not only that RSA collapses: Shor's algorithm also solves the discrete
logarithm, both over $\mathbb{Z}_p$ and over elliptic curves, so that Diffie-Hellman, DSA, ECDH, and
ECDSA collapse too, including the elliptic curves with the shortest moduli because they are immune to
index calculus.

What does not collapse is symmetric ciphers and hash functions. Grover gives a quadratic speedup in
key search, so that AES-128 drops from $2^{128}$ to about $2^{64}$ quantum operations, whereas
AES-256 stays at about $2^{128}$; the collision resistance of hashes is practically unchanged. The
answer for symmetric data is therefore simple: use keys twice as long.

Two things that are often mixed together in the news must be distinguished. Shor proves that the
factorization can be computed in polynomial time on a quantum computer without error; building the
machine is another technical matter. The estimate of Gidney and Ekera (2021) mentions about twenty
million physical qubits and eight hours for a $2048$-bit modulus, almost all of it for error
correction. That number has dropped by several orders of magnitude in one decade, but no machine has
come close. What applies now is the pattern ``harvest now, open later'': today's messages can be
stored and opened when the machine exists.

\section{Post-quantum cryptography: lattices and Kyber}
\tingkat{advanced}

\index{post-quantum cryptography}
If factorization falls, what is needed is another problem that does not fall with it. The strongest
candidates come from lattices, that is sets of points generated by integer combinations of a few
vectors.

\index{lattice}
\emph{Learning with errors} sets up the problem like this. Given a random $m\times n$ matrix $A$ over
$\mathbb{Z}_q$ and a vector $b=As+e\bmod q$, with $s$ secret and $e$ a small error, find $s$. Without
$e$ this is only a system of linear equations solved by Gaussian elimination (Chapter 8); with $e$
small but not zero, the system is almost consistent and finding $s$ is considered hard. Regev's
reduction (2005) adds to its appeal: attacking the random scheme is as hard as solving the worst case
of the lattice problem, not an average case chosen by the attacker.

Kyber, which NIST standardized in August 2024 as ML-KEM in FIPS 203, is a module variant of that idea
with $n=256$, $k=3$, and $q=3329$. Its public key is $1\,184$ bytes and its ciphertext $1\,088$
bytes, about $37$ and $34$ times the $32$-byte symmetric key it wraps, and that space cost is real.
Post-quantum signatures have a similar story: ML-DSA (Dilithium) is also lattice-based, whereas
SLH-DSA (SPHINCS+) takes the conservative hash-based road.

An honesty that is obligatory here: the lattice assumptions are much younger than the factoring
assumption. Factorization has been attacked since Gauss; LWE is only twenty years old, and the
history of cryptography is full of schemes that looked strong and then broke after ten years of
analysis. The recommended transition is therefore not replacement but hybrid: the session key is
formed from Kyber and from Diffie-Hellman exchange at once.

\section{When output becomes an oracle: padding and models}
\tingkat{advanced}

\index{padding oracle}
There is one family of attacks that does not attack the mathematics, but the way the program answers.
Bleichenbacher in 1998 showed that RSA with PKCS\#1 v1.5 padding can be opened completely if the
server answers differently for ``bad padding'' and ``good padding but wrong contents''. Every answer
leaks one bit about the message, and with a few thousand queries the message is read. The mathematics
of RSA is not violated; what is violated is the assumption that the attacker only sees the
ciphertext.

The threat pattern is the same as the one discussed in Chapter 29 and Chapter 30. A machine learning
model that returns probabilities, or error messages that differ, is an oracle; an attacker who may
ask many times uses those small differences to infer what he should not know, including whether an
example is in the training data. Chapter 30 adds a warning that applies in both worlds: a measurement
that seems to improve is not necessarily improving because the defense works.

The defensive lesson is therefore the same. Output must be made uniform: one error message for all
failures, response times that do not depend on secrets, refusals that do not distinguish the cause.
For models, this means limiting the number of queries and disguising scores. All of that reduces
leakage; it does not remove it.

\begin{tantangan}
A library uses AES-GCM with the same key $K$ for two messages, but its nonce counter is reset so that
the same $96$-bit nonce is used twice. The attacker holds both ciphers, both tags, and the contents of
the first message, namely a login page that is fixed. Show how he obtains the authentication key $H$,
then why $H$ is enough to forge a tag for a message of his own making without knowing $K$.

Hint 1: the same nonce gives the same stream key, so that the difference of the two ciphers removes
that key.

Hint 2: the tag is a polynomial in $H$ with known coefficients; two tags with the same nonce give two
equations for one $H$, and their difference becomes a polynomial whose roots can be computed.

If you get stuck, go back to the section ``Hash functions, MACs, and signatures'' and reread the
paragraph about MACs.
\end{tantangan}

\section{Worked example}
\tingkat{intermediate}

We compute one inverse by hand, because the backward steps are where mistakes happen most often. Take
$\varphi(n)=1\,048\,540$ from the toy modulus in part six and $e=65\,537$. Forward Euclid gives
\[
1\,048\,540=15\cdot65\,537+65\,485,\qquad 65\,537=1\cdot65\,485+52,
\]
\[
65\,485=1259\cdot52+17,\qquad 52=3\cdot17+1,\qquad 17=17\cdot1+0 .
\]
The last nonzero remainder is $1$, so that $\gcd=1$ and the inverse exists. Build backward starting
from the remainder one:
\[
1=52-3\cdot17=52-3(65\,485-1259\cdot52)=3778\cdot52-3\cdot65\,485 .
\]
Continue with $52=65\,537-65\,485$ and $65\,485=1\,048\,540-15\cdot65\,537$:
\[
1=3778\cdot65\,537-3781\cdot65\,485
=3778\cdot65\,537-3781(1\,048\,540-15\cdot65\,537).
\]
Group: $1=(3778+15\cdot3781)\cdot65\,537-3781\cdot1\,048\,540
=60\,493\cdot65\,537-3781\cdot1\,048\,540$. So $d=60\,493$, and the negative sign disappears
because it is read modulo $\varphi(n)$. The check is one line: $65\,537\cdot60\,493
=3\,964\,529\,741$, whereas $3781\cdot1\,048\,540=3\,964\,529\,740$; the difference is one.

Two other computations in this chapter can be checked in the same way: the trimming of
$7^{222}\bmod13$ in part three, which uses $222=12\cdot18+6$ and gives $12$, and the small key
exchange $p=23$, $g=5$, $a=4$, $b=3$ which gives $A=4$, $B=10$, and key $18$ from both directions.

\begin{resep}
Four steps to find the inverse of $a$ modulo $n$.
\begin{enumerate}
  \item Write $(a,n)$. Run Euclid: replace $(a,n)$ by $(n,a\bmod n)$ until the remainder is zero;
        record each line $n_0=q\,n_1+r$.
  \item The last nonzero remainder is $\gcd$. If it is not $1$, stop; the inverse does not exist.
  \item From the last line with remainder $1$, write $1=r-q\,s$, then replace $s$ and $r$ by their
        forms from the line above, one line up for each step.
  \item Collect the coefficient of $a$; that is the inverse. If it is negative, add $n$ until it lies
        in $0,\dots,n-1$, then check $a\cdot x\bmod n=1$.
\end{enumerate}
Warning about order: build backward from the smallest remainder upward, not from the first line.
Swapping the roles of $a$ and $n$ also changes the answer: the inverse of $a$ modulo $n$ is not the
inverse of $n$ modulo $a$.
\end{resep}

\kertas{Repeat the first computation without looking at the text. Write five lines of forward Euclid
for $65\,537$ and $1\,048\,540$, then build backward until you obtain $d$. Check by multiplying $e$
and $d$ and dividing by $\varphi(n)$: the remainder must be one. After that, answer: why can the same
steps not be used to open RSA without the secret key, even though this whole computation needs only
five divisions?}

\section{Computational lab}
\tingkat{advanced}

The following program uses only the standard library and runs in under one second. It checks four
claims in the body of the chapter: the units modulo $26$, the statistics of Euclid's steps, the
equality of the Diffie-Hellman keys, and the looseness of plain RSA. The last part tests the birthday
prediction on SHA-256 truncated to two bytes.

\begin{lab}
\begin{verbatim}
import hashlib, math, random

def egcd(a, b):
    if b == 0:
        return a, 1, 0
    g, x, y = egcd(b, a % b)
    return g, y, x - (a // b) * y

def invers(a, n):
    g, x, _ = egcd(a, n)
    return None if g != 1 else x % n

def langkah(a, b):
    k = 0
    while b:
        a, b = b, a % b
        k += 1
    return k

def orde(a, n):
    k, y = 1, a % n
    while y != 1:
        y = y * a % n
        k += 1
    return k

unit = [a for a in range(26) if invers(a, 26) is not None]
print("units mod 26:", len(unit), unit)
print("inverse 3 mod 26 =", invers(3, 26), " inverse 2 mod 26 =", invers(2, 26))

g0, x0, y0 = egcd(1071, 462)
print("gcd(1071,462) =", g0, " x,y =", x0, y0, " steps =", langkah(1071, 462))

p, g, a, b = 10007, 5, 1234, 5678
A, B = pow(g, a, p), pow(g, b, p)
print("DH: A =", A, " B =", B, " K =", pow(B, a, p), pow(A, b, p))
print("order(5) =", orde(5, p), " order(5^5003) =", orde(pow(5, 5003, p), p))

random.seed(29)
L = [langkah(random.getrandbits(32), random.getrandbits(32))
     for _ in range(20000)]
print("euclid 32-bit: mean %.3f, max %d, min %d" % (sum(L) / len(L), max(L), min(L)))

p, q, e = 1019, 1031, 3
n, phi = p * q, (p - 1) * (q - 1)
d = invers(e, phi)
print("RSA: n =", n, " bits =", n.bit_length(), " phi =", phi, " d =", d)
for m in (424242, 42):
    c = pow(m, e, n)
    print("  m =", m, " c =", c, " D(c) =", pow(c, d, n), " m^3 =", m ** 3)
kali = pow(1234, e, n) * pow(567, e, n) % n
print("E(1234)E(567) mod n =", kali, " D(kali) =", pow(kali, d, n),
      " 1234*567 mod n =", 1234 * 567 % n)

def tabrakan():
    terlihat, k = {}, 0
    while True:
        m = random.getrandbits(32)
        h = hashlib.sha256(m.to_bytes(4, "big")).digest()[:2]
        k += 1
        if h in terlihat:
            return k
        terlihat[h] = m

T = [tabrakan() for _ in range(20)]
print("16-bit collisions: mean %.1f, min %d, max %d" % (sum(T) / 20, min(T), max(T)))
print("prediction sqrt(pi/2 * 2^16) =", round(math.sqrt(math.pi / 2 * 2**16), 1))
\end{verbatim}

The output, as it stands:

\begin{verbatim}
units mod 26: 12 [1, 3, 5, 7, 9, 11, 15, 17, 19, 21, 23, 25]
inverse 3 mod 26 = 9  inverse 2 mod 26 = None
gcd(1071,462) = 21  x,y = -3 7  steps = 3
DH: A = 4873  B = 3189  K = 5674 5674
order(5) = 10006  order(5^5003) = 2
euclid 32-bit: mean 18.750, max 32, min 6
RSA: n = 1050589  bits = 21  phi = 1048540  d = 699027
  m = 424242  c = 618000  D(c) = 424242  m^3 = 76355615883580488
  m = 42  c = 74088  D(c) = 42  m^3 = 74088
E(1234)E(567) mod n = 551173  D(kali) = 699678  1234*567 mod n = 699678
16-bit collisions: mean 362.7, min 100, max 743
prediction sqrt(pi/2 * 2^16) = 320.8
\end{verbatim}
\end{lab}

The first line confirms part one: twelve units modulo $26$, with $1$ and $25$ their own inverses and
the remaining ten paired; that number $12$ is the same as $\varphi(26)$.

The second and third lines show two faces of Euclid: three divisions for $1071$ and $462$ with Bezout
coefficients $(-3,7)$, and an average of $18.75$ divisions with maximum $32$ for two thousand random
$32$-bit pairs. The fourth and fifth lines are the promised key exchange: $4873$ and $3189$ on the
channel, $5674$ at both ends, whereas $5^{5003}$ has order $2$ and is therefore the trap of part five.

The three RSA lines show two opposite things. For $m=424\,242$ everything works: the cipher
$618\,000$ returns to $424\,242$. For $m=42$ with $e=3$, the cipher is $74\,088$, exactly $42^3$, so
that no key needs to be found, only a cube root; the next line shows the multiplicative property: the
product of two ciphers opens the product message.

The last four lines test the birthday prediction: an average of $362.7$ samples against a prediction
of $320.8$, with the wide spread expected from an exponential distribution.

\begin{untukdosen}
The teaching order I suggest: start from the student's cipher in the opening, then the
well-definedness lemma and the unit theorem in the first meeting; leave Diffie-Hellman and RSA for the
second meeting, and put Shor and post-quantum as reading for the third. Do not teach Euclid before
asking ``when does division exist''.

Two misconceptions almost always appear. First, students think $\varphi(n)=n-1$ for all $n$, then are
startled when $\varphi(26)=12$; cure it by asking them to list the units modulo $26$ on the board
before any formula is mentioned. Second, they think RSA is secure because $n$ is hard to factor,
without mentioning the size of $n$ and without mentioning that this difficulty is an assumption; ask
them to write the security claim in the form ``if ... then ...''.

A blackboard activity: two groups compute $3^{12}\bmod26$, one by repeated multiplication and one by
Fermat's little theorem, then compare time and results. A sign of readiness: they can explain without
notes why $[2]_{26}$ has no inverse.
\end{untukdosen}

\section{Summary}
\begin{ringkasan}
\begin{itemize}
    \item Arithmetic modulo $n$ works on residue classes, and division exists only for classes that
        share no factor with $n$; Euclid gives the inverse in time proportional to the number of
        digits.
    \item The shuffling lemma, namely that multiplication by a unit is only a shuffling of the order,
        gives Euler's theorem in five lines and Fermat's little theorem as a corollary. The order
        divides $\varphi(n)$, and a primitive root is an element of full order.
    \item Diffie-Hellman agrees on a key through the equality $g^{ab}$ from two directions, with
        security conditional on the difficulty of the discrete logarithm; without authentication it
        falls to the man-in-the-middle attack.
    \item RSA reverses the direction: the public key locks, the secret key opens, and the proof uses
        Fermat's little theorem on both prime factors. Plain RSA is deterministic, multiplicative, and
        does not protect small messages; padding closes all three.
    \item A hash function gives a fingerprint, a MAC gives authenticity between two parties who share
        a key, a signature gives evidence that anyone can check; collision resistance is only half the
        output length, and the numbers $362.7$ against $320.8$ in the lab show that the prediction
        works.
    \item Every guarantee in this chapter is conditional. Shor brings down RSA and the whole of
        discrete logarithms if a large-scale quantum computer exists; symmetric ciphers and hashes
        only need to double the key length, and lattice schemes such as ML-KEM are the possible
        replacements.
\end{itemize}
\end{ringkasan}

\section{Exercises}
\begin{latihan}
\begin{enumerate}
  \item Assess this claim: ``Because $2^{100}\equiv1\pmod{101}$, the number $101$ is proved prime.''
        Decide true, false, or not enough evidence, and name what that congruence does prove.
  \item Prove that the product of two units modulo $n$ is also a unit, then conclude that $U(n)$ is
        closed under multiplication; show also that if $a$ is not a unit, there is $b\neq0$ with
        $ab\equiv0\pmod n$.
  \item Prove that for prime $p$, $a^2\equiv1\pmod p$ is satisfied only by $a\equiv1$ or
        $a\equiv-1$; explain why the same fails for modulus $8$, and name all square roots of one
        modulo $8$.
  \item Prove $\varphi(pq)=\varphi(p)\varphi(q)$ for distinct primes $p,q$ by counting the classes
        $0\le a<pq$ that share a factor with $pq$, then check it on $p=1019$ and $q=1031$.
  \item Compute $17^{-1}\bmod3120$ with backward Euclid, then check it with Python. Repeat for
        $65\,537^{-1}\bmod1\,048\,540$ and compare the number of divisions of the two.
  \item Compute $5^{100}\bmod101$, $3^{200}\bmod101$, and $2^{1000}\bmod101$ using Fermat's little
        theorem, without computing the full power. Check all three with \texttt{pow(a,k,m)}.
  \item On modulus $23$ with base $5$, Alice uses $a=7$ and Bob $b=15$. Compute $A$, $B$, and the
        shared key, then compute the key obtained by an active eavesdropper who replaces $A$ with
        $11$ in front of Bob, and explain why two different keys appear.
  \item Explain why RSA with $e=3$ and a message $m<n^{1/3}$ can be opened without the secret key,
        how long the longest message affected on the modulus $1\,050\,589$ is, and what padding does
        to prevent it.
  \item Assess this claim: ``ML-KEM is already proved secure, so we can stop using RSA now.'' Name
        two things that must be true for the first sentence to make sense, and one reason why the
        second sentence remains wrong even if the first is true.
  \item A system uses a hash with a $48$-bit output. How many samples are needed for a collision
        probability of one in two, and how many if the output is truncated to $32$ bits?
\end{enumerate}

Hints for the three hardest items. Item 3: write $a^2-1=(a-1)(a+1)$ and use Euclid's lemma; for
modulus $8$, find the roots by trying $a=0,\dots,7$. Item 7: compute $11^7\bmod23$ as the key that
Bob sees, then compare it with $11^{15}\bmod23$ computed by Alice. Item 9: distinguish ``proved
secure'' from ``secure against known attacks''.
\end{latihan}

\begin{untukpeneliti}
What is really proved in this chapter is its mathematical part: Bezout's theorem, the unit theorem,
Euler's theorem, the order lemma, and the RSA theorem all have complete proofs and do not depend on
any conjecture. What is only popularized is the word ``secure''. No theorem states that RSA-2048 is
hard; what exists are cost estimates for the best known algorithms, and those estimates change
whenever cryptanalysis makes progress.

The primary sources: Diffie and Hellman (1976), sections 3--4; Rivest, Shamir, and Adleman (1978),
sections 2--4; Shor (1997), section 5 for the reduction of factorization to order finding; Regev
(2005), section 3 for the LWE reduction; Bleichenbacher (1998); Gidney and Ekera (2021); FIPS 203 and
204.

Limits of validity that must be mentioned: security reductions usually produce loose bounds, the
attack models do not contain side channels, and padding attacks show that the difference of one error
message is enough to bring down a scheme that is proved secure in the model. Open questions: how tight
the reduction for ML-KEM is, whether there is a faster quantum algorithm for lattice problems, and how
to prove the security of schemes used together with secret data without leaking access patterns.
\end{untukpeneliti}

\peta{Chapter 31 follows Diffie and Hellman (1976), section 3 for key exchange; Rivest, Shamir, and
Adleman (1978), sections 2--4 for RSA and signatures; Shor (1997), sections 5--6 for the reduction of
factorization to order finding; Regev (2005), section 3 for the LWE reduction; Bleichenbacher (1998);
Gidney and Ekera (2021); FIPS 203 and 204. On Chapter 2 for bijections and partitions, Chapter 3 for
the Archimedean property and the discrete form of the infimum property, Chapter 8 for Gaussian
elimination, Chapter 26 for the reason a security claim cannot be a theorem without conditions, and
Chapters 29--30 for the oracle threat pattern.}
